\documentclass[12pt,a4paper]{article}
\usepackage[utf8]{inputenc}
\usepackage[T1]{fontenc}
\usepackage{geometry}
\geometry{a4paper, margin=1in, top=0.8in, bottom=0.8in}

% Packages for design and math
\usepackage{amsmath, amssymb, amsthm}
\usepackage{mathtools}
\usepackage{physics}
\usepackage{bm}
\usepackage{graphicx}
\usepackage{xcolor}
\usepackage{tcolorbox}
\usepackage{tikz}
\usepackage{enumitem}
\usepackage{fancyhdr}
\usepackage{hyperref}
\usepackage{listings}
\usepackage{booktabs}
\usepackage{array}
\usepackage{colortbl}
\usepackage{multirow}
\usepackage{pgfplots}
\pgfplotsset{compat=1.18}

% Color definitions
\definecolor{primary}{RGB}{25, 55, 95}
\definecolor{secondary}{RGB}{70, 130, 200}
\definecolor{accent}{RGB}{220, 70, 50}
\definecolor{formula}{RGB}{200, 80, 60}
\definecolor{mathematica}{RGB}{50, 100, 80}
\definecolor{lightbg}{RGB}{240, 248, 255}

% Custom tcolorbox styles
\newtcolorbox{definitionbox}[1][]{
    colback=lightbg,
    colframe=primary,
    coltitle=white,
    fonttitle=\bfseries\large,
    title={#1},
    sharp corners,
    boxrule=1.5pt,
    left=10pt,
    right=10pt,
    top=8pt,
    bottom=8pt
}

\newtcolorbox{formulabox}[1][]{
    colback=yellow!15,
    colframe=formula,
    sharp corners,
    boxrule=1pt,
    left=8pt,
    right=8pt,
    top=5pt,
    bottom=5pt,
    #1
}

\newtcolorbox{mathematicabox}{
    colback=gray!10,
    colframe=mathematica,
    sharp corners,
    boxrule=1pt,
    fonttitle=\bfseries,
    title=Wolfram Mathematica 13 Implementation,
    left=8pt,
    right=8pt,
    top=5pt,
    bottom=5pt
}

\newtcolorbox{tipbox}{
    colback=green!10,
    colframe=green!50!black,
    sharp corners,
    boxrule=1pt,
    left=8pt,
    right=8pt,
    top=5pt,
    bottom=5pt
}

% Header and footer
\pagestyle{fancy}
\fancyhf{}
\fancyhead[L]{\textcolor{primary}{\textbf{Linear Algebra \& Calculus}}}
\fancyhead[R]{\textcolor{primary}{\textbf{Exam Preparation Notes}}}
\fancyfoot[C]{\thepage}
\fancyfoot[R]{\textcolor{gray}{Wolfram Mathematica 13 Lab}}

% Section formatting
\usepackage{titlesec}
\titleformat{\section}
    {\color{primary}\normalfont\Large\bfseries}
    {\color{primary}\thesection}{1em}{}
\titleformat{\subsection}
    {\color{secondary}\normalfont\large\bfseries}
    {\color{secondary}\thesubsection}{1em}{}

% Math operators
\providecommand{\rank}{\operatorname{rank}}
\providecommand{\nullity}{\operatorname{nullity}}
\providecommand{\proj}{\operatorname{proj}}
\providecommand{\comp}{\operatorname{comp}}
\providecommand{\tr}{\operatorname{tr}}

% Document
\begin{document}

% Title Page
\begin{titlepage}
    \centering
    \vspace*{2cm}
    
    {\Huge\bfseries\color{primary} Complete Exam Preparation Notes}
    
    \vspace{0.5cm}
    {\Large\bfseries\color{secondary} Linear Algebra \& Calculus}
    
    \vspace{0.5cm}
    {\large with Wolfram Mathematica 13 Implementation}
    
    \vspace{2cm}
    
    \begin{tcolorbox}[colback=lightbg, colframe=primary, sharp corners, width=0.8\textwidth, center]
        \centering
        \textbf{Covering Assignments:}\\
        \vspace{0.3cm}
        1B — Complex Numbers \& Series \\
        2B — Definite Integrals \& Area \\
        3B — Lines, Planes \& Vectors \\
        4A — Matrix Operations \& Linear Systems \\
        4B — Linear Independence \& Eigenvalues
    \end{tcolorbox}
    
    \vspace{2cm}
    
    \textcolor{gray}{Department of Applied Mathematics}\\
    \textcolor{gray}{University of Dhaka}
    
    \vfill
    \textcolor{gray}{\today}
\end{titlepage}

% Table of Contents
\tableofcontents
\newpage

\section*{How to Use This Note (Wolfram Mathematica 13)}
\addcontentsline{toc}{section}{How to Use This Note (Wolfram Mathematica 13)}

\begin{definitionbox}[Recommended Study Flow]
    Each topic is organized in a practical exam-ready sequence:
    \begin{enumerate}
        \item \textbf{Core explanation:} understand the concept and conditions.
        \item \textbf{Formula focus:} memorize the key equation(s).
        \item \textbf{Worked example:} follow the solved model problem.
        \item \textbf{Mathematica 13 code:} verify and automate the calculation.
        \item \textbf{Quick review:} use the Master Formula Sheet and tips section.
    \end{enumerate}
\end{definitionbox}

\begin{tipbox}
    \textbf{Important for Mathematica 13:}
    \begin{itemize}
        \item Use exact arithmetic first (fractions/symbolic forms), then apply \texttt{N[...]}.
        \item For linear algebra, use built-ins such as \texttt{Det}, \texttt{Inverse}, \texttt{LinearSolve}, and \texttt{Eigensystem}.
        \item Always cross-check symbolic output with a numerical or graphical check when possible.
    \end{itemize}
\end{tipbox}
\newpage

% ============================================================
% SECTION 1: ASSIGNMENT 1B - COMPLEX NUMBERS & SERIES
% ============================================================
\section{Assignment 1B: Complex Numbers \& Series}

\begin{definitionbox}[Learning Goals]
    Convert complex inequalities to geometric regions, estimate possible real roots using Descartes' rule,
    evaluate finite/infinite-style series patterns, compare AM--GM--HM correctly, and implement all methods in Mathematica 13.
\end{definitionbox}

\subsection{Complex Inequalities}

\begin{definitionbox}[Key Concept]
    For a complex number \(z = x + iy\), the modulus is defined as:
    \[ |z| = \sqrt{x^2 + y^2} \]
    
    To convert a complex inequality to real form:
    \begin{enumerate}
        \item Substitute \(z = x + iy\)
        \item Use \( \left| \frac{z - a}{z - b} \right| = \frac{\sqrt{(x-a)^2 + y^2}}{\sqrt{(x-b)^2 + y^2}} \)
        \item Square both sides (both sides positive)
        \item Simplify to get circle/region equation
    \end{enumerate}
\end{definitionbox}

\begin{formulabox}
    \textbf{Example:} \( \left| \frac{z-3}{z+3} \right| < 2 \)
    
    \[
    \frac{\sqrt{(x-3)^2 + y^2}}{\sqrt{(x+3)^2 + y^2}} < 2
    \implies (x+5)^2 + y^2 > 16
    \]
    
    \textbf{Result:} Region outside circle centered at \((-5,0)\) with radius \(4\).
\end{formulabox}

\begin{mathematicabox}
    \begin{lstlisting}[language=Mathematica, basicstyle=\ttfamily\small]
(* Solve complex inequality *)
ineq = Abs[(z - 3)/(z + 3)] < 2;
solution = Reduce[ineq, z, Complexes];
Simplify[solution]

(* Plot the region *)
RegionPlot[(x + 5)^2 + y^2 > 16, {x, -12, 2}, {y, -6, 6},
  PlotLabel -> "|(z-3)/(z+3)| < 2",
  AxesLabel -> {"Re(z)", "Im(z)"}]
    \end{lstlisting}
\end{mathematicabox}

\subsection{Descartes' Rule of Signs}

\begin{definitionbox}[Descartes' Rule of Signs]
    For a polynomial \(f(x)\) with real coefficients:
    \begin{itemize}
        \item \textbf{Positive real roots:} Number equals sign changes in \(f(x)\) (or less by an even number)
        \item \textbf{Negative real roots:} Number equals sign changes in \(f(-x)\) (or less by an even number)
    \end{itemize}
\end{definitionbox}

\begin{formulabox}
    \textbf{Example:} \(f(x) = x^9 + 5x^8 - x^3 + 7x + 2\)
    
    \begin{tabular}{ll}
    \(f(x)\) signs: & \(+, +, -, +, +\) → 2 sign changes \\
    \(f(-x)\) signs: & \(-, +, +, -, +\) → 3 sign changes \\
    \hline
    Positive real roots: & 0 or 2 \\
    Negative real roots: & 1 or 3 \\
    \end{tabular}
\end{formulabox}

\begin{mathematicabox}
    \begin{lstlisting}[language=Mathematica, basicstyle=\ttfamily\small]
p[x_] := x^9 + 5x^8 - x^3 + 7x + 2;

(* Count sign changes *)
coeff = CoefficientList[p[x], x];
signChanges = Count[Differences[Sign[coeff]], 2] + 
              Count[Differences[Sign[coeff]], -2];

(* Find actual roots *)
roots = NSolve[p[x] == 0, x];
realRoots = Select[roots, Im[#[[1,2]]] == 0 &];
    \end{lstlisting}
\end{mathematicabox}

\subsection{Series Summation Using Loops}

\begin{definitionbox}[Loop Structures in Mathematica]
    \begin{itemize}
        \item \textbf{Do loop:} \texttt{Do[expression, \{i, imin, imax\}]}
        \item \textbf{For loop:} \texttt{For[start, test, increment, body]}
        \item \textbf{While loop:} \texttt{While[condition, body]}
    \end{itemize}
\end{definitionbox}

\begin{formulabox}
    \textbf{Series:} \( \displaystyle \frac{1}{3\cdot4\cdot5} + \frac{1}{4\cdot5\cdot6} + \frac{1}{5\cdot6\cdot7} + \cdots \)
    
    General term: \( a_k = \frac{1}{k(k+1)(k+2)} \), starting at \(k=3\)
    
    Telescoping form: \( a_k = \frac{1}{2}\left[\frac{1}{k(k+1)} - \frac{1}{(k+1)(k+2)}\right] \)
\end{formulabox}

\begin{mathematicabox}
    \begin{lstlisting}[language=Mathematica, basicstyle=\ttfamily\small]
(* Do loop *)
sum = 0;
Do[sum += 1/(k*(k+1)*(k+2)), {k, 3, 102}]

(* For loop *)
sum = 0;
For[k = 3, k <= 102, k++, 
  sum += 1/(k*(k+1)*(k+2))]

(* Built-in Sum *)
Sum[1/(k*(k+1)*(k+2)), {k, 3, 102}]

(* Telescoping sum exact value *)
exact = 1/2*(1/(3*4) - 1/(103*104))
    \end{lstlisting}
\end{mathematicabox}

\subsection{Means: AM, GM, HM}

\begin{definitionbox}[Three Classical Means]
    For positive numbers \(x_1, x_2, \ldots, x_n\):
    
    \begin{align*}
    \text{Arithmetic Mean (AM)} &= \frac{1}{n}\sum_{i=1}^n x_i \\
    \text{Geometric Mean (GM)} &= \left(\prod_{i=1}^n x_i\right)^{1/n} \\
    \text{Harmonic Mean (HM)} &= \frac{n}{\sum_{i=1}^n \frac{1}{x_i}}
    \end{align*}
    
    \textbf{AM–GM Inequality:} \(AM \geq GM \geq HM\) (with equality iff all numbers equal)
\end{definitionbox}

\begin{mathematicabox}
    \begin{lstlisting}[language=Mathematica, basicstyle=\ttfamily\small]
(* Generate 15 random integers *)
randomList = RandomInteger[{1, 20}, 15];

(* Calculate means *)
AM = Mean[randomList];
GM = (Product[randomList[[i]], {i,1,15}])^(1/15);
HM = 15 / Sum[1/randomList[[i]], {i,1,15}];

(* Check inequality *)
AM >= GM && GM >= HM
    \end{lstlisting}
\end{mathematicabox}

\subsection{Voting Theory Methods}

\begin{definitionbox}[Voting Methods]
    \begin{tabular}{|p{3cm}|p{10cm}|}
        \hline
        \textbf{Method} & \textbf{Description} \\
        \hline
        Borda Count & Points assigned by rank (4 for 1st, 3 for 2nd, 2 for 3rd, 1 for 4th) \\
        Plurality & Candidate with most first-choice votes wins \\
        Plurality-with-Elimination & Eliminate lowest, redistribute votes \\
        Pairwise Comparison & Each pair compared head-to-head \\
        \hline
    \end{tabular}
\end{definitionbox}

\begin{formulabox}
    \textbf{Mini Example (Borda Count):} Suppose 3 candidates \(A,B,C\) and 10 voters.
    If first/second/third place points are \(3,2,1\), then total score is
    \[\text{Score}(X)=3\cdot n_1(X)+2\cdot n_2(X)+1\cdot n_3(X)\]
    Candidate with highest score wins.
\end{formulabox}

\begin{mathematicabox}
    \begin{lstlisting}[language=Mathematica, basicstyle=\ttfamily\small]
(* Borda Count demo: columns = rank positions for each voter block *)
candidates = {"A", "B", "C"};
points = {3, 2, 1};

(* Example preference blocks and voter counts *)
ballots = {{"A", "B", "C"}, {"B", "C", "A"}, {"C", "A", "B"}};
counts = {4, 3, 3};

scores = AssociationThread[candidates -> ConstantArray[0, Length[candidates]]];
Do[
  Do[
    scores[ballots[[j, r]]] += points[[r]]*counts[[j]],
    {r, 1, Length[points]}
  ],
  {j, 1, Length[ballots]}
];

scores
    \end{lstlisting}
\end{mathematicabox}

% ============================================================
% SECTION 2: ASSIGNMENT 2B - INTEGRALS & AREA
% ============================================================
\section{Assignment 2B: Definite Integrals \& Area}

\begin{definitionbox}[Learning Goals]
    Compute definite integrals geometrically and analytically, use Riemann sums, find area between curves,
    evaluate average values, and solve motion/volume/arc-length problems with Mathematica 13.
\end{definitionbox}

\subsection{Definite Integrals and Geometric Areas}

\begin{definitionbox}[Fundamental Theorem of Calculus]
    \[ \int_a^b f(x)\,dx = F(b) - F(a) \quad \text{where } F'(x) = f(x) \]
\end{definitionbox}

\begin{formulabox}
    \textbf{Common Geometric Areas:}
    \begin{align*}
    \text{Rectangle:} &\quad \int_a^b c\,dx = c(b-a) \\
    \text{Triangle:} &\quad \int_a^b (mx + c)\,dx = \frac{1}{2}(b-a)(f(a)+f(b)) \\
    \text{Circle:} &\quad \int_0^r \sqrt{r^2 - x^2}\,dx = \frac{\pi r^2}{4}
    \end{align*}
\end{formulabox}

\begin{mathematicabox}
    \begin{lstlisting}[language=Mathematica, basicstyle=\ttfamily\small]
(* Definite integrals *)
Integrate[2, {x, 1, 4}]           (* Rectangle: 3*2=6 *)
Integrate[x+2, {x, -1, 2}]        (* Trapezoid *)
Integrate[Sqrt[1-x^2], {x, 0, 1}] (* Quarter circle: pi/4 *)

(* Numerical approximation *)
NIntegrate[Sqrt[1-x^2], {x, 0, 1}]
    \end{lstlisting}
\end{mathematicabox}

\subsection{Riemann Sums}

\begin{definitionbox}[Riemann Sum]
    \[ \int_a^b f(x)\,dx \approx \sum_{i=1}^n f(x_i^*) \Delta x \]
    
    where \(\Delta x = \frac{b-a}{n}\) and:
    \begin{itemize}
        \item Left endpoint: \(x_i^* = a + (i-1)\Delta x\)
        \item Right endpoint: \(x_i^* = a + i\Delta x\)
        \item Midpoint: \(x_i^* = a + (i-\frac{1}{2})\Delta x\)
    \end{itemize}
\end{definitionbox}

\begin{mathematicabox}
    \begin{lstlisting}[language=Mathematica, basicstyle=\ttfamily\small]
f[x_] := 9 - x^2;
a = 0; b = 3; n = 10;
deltaX = (b - a)/n;

(* Left sum *)
leftSum = Sum[f[a + (i-1)*deltaX]*deltaX, {i, 1, n}];

(* Right sum *)
rightSum = Sum[f[a + i*deltaX]*deltaX, {i, 1, n}];

(* Midpoint sum *)
midSum = Sum[f[a + (i-0.5)*deltaX]*deltaX, {i, 1, n}];

(* Exact value *)
exact = Integrate[f[x], {x, a, b}];
    \end{lstlisting}
\end{mathematicabox}

\subsection{Area Between Curves}

\begin{formulabox}
    \[ \text{Area} = \int_a^b |f(x) - g(x)|\,dx \]
    
    where \(a\) and \(b\) are intersection points of \(f\) and \(g\).
\end{formulabox}

\begin{mathematicabox}
    \begin{lstlisting}[language=Mathematica, basicstyle=\ttfamily\small]
(* Area between y=x^2 and y=2x+4 *)
f[x_] := x^2;
g[x_] := 2x + 4;

intersections = Solve[f[x] == g[x], x];
area = Integrate[g[x] - f[x], {x, -2, 4}]
(* Result: 36 *)

Plot[{f[x], g[x]}, {x, -3, 5}, Filling -> {1 -> {2}}]
    \end{lstlisting}
\end{mathematicabox}

\subsection{Average Value of a Function}

\begin{formulabox}
    \[ f_{\text{avg}} = \frac{1}{b-a} \int_a^b f(x)\,dx \]
\end{formulabox}

\begin{mathematicabox}
    \begin{lstlisting}[language=Mathematica, basicstyle=\ttfamily\small]
(* Drug concentration problem *)
C[t_] := (0.25 t)/(t^2 + 4);
time = 3 + 17/60;  (* 3 hours 17 minutes *)

avgConcentration = (1/time) * Integrate[C[t], {t, 0, time}];
N[avgConcentration]
    \end{lstlisting}
\end{mathematicabox}

\subsection{Free Fall Motion}

\begin{formulabox}
    \begin{align*}
    s(t) &= s_0 + v_0 t - \frac{1}{2}gt^2 \\
    v(t) &= v_0 - gt
    \end{align*}
    
    where \(g = 32\ \text{ft/s}^2\) or \(9.8\ \text{m/s}^2\).
\end{formulabox}

\begin{mathematicabox}
    \begin{lstlisting}[language=Mathematica, basicstyle=\ttfamily\small]
(* Object dropped from 1276 ft *)
s[t_] := 1276 - 16t^2;
timeToGround = Solve[s[t] == 0, t]  (* t = Sqrt[1276/16] *)
impactVelocity = -32 * timeToGround[[2,1,2]]
    \end{lstlisting}
\end{mathematicabox}

\subsection{Volume of Solids of Revolution}

\begin{definitionbox}[Volume Methods]
    \begin{tabular}{|p{4cm}|p{10cm}|}
        \hline
        \textbf{Method} & \textbf{Formula} \\
        \hline
        Disk (x-axis) & \(V = \pi \int_a^b [R(x)]^2\,dx\) \\
        Washer & \(V = \pi \int_a^b ([R(x)]^2 - [r(x)]^2)\,dx\) \\
        Shell (y-axis) & \(V = 2\pi \int_a^b x f(x)\,dx\) \\
        \hline
    \end{tabular}
\end{definitionbox}

\begin{mathematicabox}
    \begin{lstlisting}[language=Mathematica, basicstyle=\ttfamily\small]
(* Volume: y=Sqrt[x] from x=1 to 6 about x-axis *)
f[x_] := Sqrt[x];
volume = pi * Integrate[f[x]^2, {x, 1, 6}]
(* = 35pi/2 *)

(* Shell method: y=x^2 from 0 to 2 about y=-1 *)
volumeShell = 2pi * Integrate[(x^2+1)^2 - 1^2, {x, 0, 2}]
    \end{lstlisting}
\end{mathematicabox}

\subsection{Arc Length and Surface Area}

\begin{formulabox}
    \textbf{Arc Length:}
    \[ L = \int_a^b \sqrt{1 + \left(\frac{dy}{dx}\right)^2}\,dx \]
    
    \textbf{Surface Area (about x-axis):}
    \[ S = 2\pi \int_a^b f(x) \sqrt{1 + [f'(x)]^2}\,dx \]
    
    \textbf{Parametric Arc Length:}
    \[ L = \int_a^b \sqrt{\left(\frac{dx}{dt}\right)^2 + \left(\frac{dy}{dt}\right)^2}\,dt \]
\end{formulabox}

% ============================================================
% SECTION 3: ASSIGNMENT 3B - LINES, PLANES & VECTORS
% ============================================================
\section{Assignment 3B: Lines, Planes \& Vectors}

\begin{definitionbox}[Learning Goals]
    Write equations of lines and planes, compute distances and angles, apply dot/cross products,
    and solve 3D geometry questions efficiently in Mathematica 13.
\end{definitionbox}

\subsection{Lines in 3D Space}

\begin{definitionbox}[Parametric Equation of a Line]
    \[ \mathbf{r}(t) = \mathbf{r}_0 + t\mathbf{v} \]
    
    where \(\mathbf{r}_0\) is a point on the line and \(\mathbf{v}\) is the direction vector.
\end{definitionbox}

\begin{formulabox}
    \textbf{Distance from point to line:}
    \[ d = \frac{|(\mathbf{P} - \mathbf{Q}) \times \mathbf{v}|}{|\mathbf{v}|} \]
    
    \textbf{Angle between lines:}
    \[ \cos\theta = \frac{|\mathbf{v}_1 \cdot \mathbf{v}_2|}{|\mathbf{v}_1||\mathbf{v}_2|} \]
    
    \textbf{Skew lines distance:}
    \[ d = \frac{|(\mathbf{P}_2 - \mathbf{P}_1) \cdot (\mathbf{v}_1 \times \mathbf{v}_2)|}{|\mathbf{v}_1 \times \mathbf{v}_2|} \]
\end{formulabox}

\begin{mathematicabox}
    \begin{lstlisting}[language=Mathematica, basicstyle=\ttfamily\small]
(* Parametric line *)
line = p0 + t*v

(* Distance point to line *)
distance = Norm[Cross[P-Q, v]] / Norm[v]

(* Angle between lines *)
cosTheta = Abs[v1.v2]/(Norm[v1]*Norm[v2])

(* Skew lines distance *)
dist = Abs[(p2-p1).Cross[v1,v2]] / Norm[Cross[v1,v2]]
    \end{lstlisting}
\end{mathematicabox}

\subsection{Planes in 3D Space}

\begin{definitionbox}[Equation of a Plane]
    \[ \mathbf{n} \cdot (\mathbf{r} - \mathbf{r}_0) = 0 \quad \text{or} \quad Ax + By + Cz + D = 0 \]
    
    where \(\mathbf{n} = (A,B,C)\) is the normal vector.
\end{definitionbox}

\begin{formulabox}
    \textbf{Distance from point to plane:}
    \[ d = \frac{|Ax_0 + By_0 + Cz_0 + D|}{\sqrt{A^2 + B^2 + C^2}} \]
    
    \textbf{Angle between planes:}
    \[ \cos\theta = \frac{|\mathbf{n}_1 \cdot \mathbf{n}_2|}{|\mathbf{n}_1||\mathbf{n}_2|} \]
    
    \textbf{Plane from 3 points:}
    \[ \mathbf{n} = (\mathbf{Q}-\mathbf{P}) \times (\mathbf{R}-\mathbf{P}) \]
\end{formulabox}

\begin{mathematicabox}
    \begin{lstlisting}[language=Mathematica, basicstyle=\ttfamily\small]
(* Normal from 3 points *)
n = Cross[Q-P, R-P];

(* Plane equation *)
planeEq = n.({x,y,z} - P) == 0

(* Distance point to plane *)
dist = Abs[A*x0+B*y0+C*z0+D]/Sqrt[A^2+B^2+C^2]
    \end{lstlisting}
\end{mathematicabox}

\subsection{Vector Operations}

\begin{definitionbox}[Vector Products]
    \begin{align*}
    \text{Dot product:} &\quad \mathbf{a} \cdot \mathbf{b} = a_1b_1 + a_2b_2 + a_3b_3 \\
    \text{Cross product:} &\quad \mathbf{a} \times \mathbf{b} = \begin{vmatrix}
    \mathbf{i} & \mathbf{j} & \mathbf{k} \\
    a_1 & a_2 & a_3 \\
    b_1 & b_2 & b_3
    \end{vmatrix} \\
    \text{Scalar triple product:} &\quad \mathbf{a} \cdot (\mathbf{b} \times \mathbf{c})
    \end{align*}
\end{definitionbox}

\begin{formulabox}
    \textbf{Applications:}
    \begin{itemize}
        \item Area of triangle: \(\frac{1}{2}|\mathbf{AB} \times \mathbf{AC}|\)
        \item Volume of parallelepiped: \(|\mathbf{a} \cdot (\mathbf{b} \times \mathbf{c})|\)
        \item Work: \(W = \mathbf{F} \cdot \mathbf{d}\)
        \item Torque: \(\boldsymbol{\tau} = \mathbf{r} \times \mathbf{F}\)
    \end{itemize}
\end{formulabox}

\begin{mathematicabox}
    \begin{lstlisting}[language=Mathematica, basicstyle=\ttfamily\small]
(* Vector operations *)
a = {a1, a2, a3};
b = {b1, b2, b3};
dot = a.b
cross = Cross[a, b]
mag = Norm[a]
unit = a/Norm[a]

(* Area of triangle *)
area = Norm[Cross[B-A, C-A]]/2

(* Volume of parallelepiped *)
vol = Abs[Dot[a, Cross[b, c]]]
    \end{lstlisting}
\end{mathematicabox}

% ============================================================
% SECTION 4: ASSIGNMENT 4A - MATRIX OPERATIONS
% ============================================================
\section{Assignment 4A: Matrix Operations \& Linear Systems}

\begin{definitionbox}[Learning Goals]
    Perform matrix operations correctly, build special matrices, solve linear systems,
    and interpret applied models such as network flow using Mathematica 13 commands.
\end{definitionbox}

\subsection{Matrix Operations}

\begin{definitionbox}[Basic Matrix Operations]
    \begin{tabular}{|p{3cm}|p{10cm}|}
        \hline
        \textbf{Operation} & \textbf{Mathematica Syntax} \\
        \hline
        Addition & \texttt{A + B} \\
        Subtraction & \texttt{A - B} \\
        Scalar multiplication & \texttt{k * A} \\
        Matrix multiplication & \texttt{A . B} (dot, not *) \\
        Transpose & \texttt{Transpose[A]} \\
        Inverse & \texttt{Inverse[A]} \\
        Determinant & \texttt{Det[A]} \\
        Trace & \texttt{Tr[A]} \\
        \hline
    \end{tabular}
\end{definitionbox}

\begin{tipbox}
    \textbf{Important:} Matrix multiplication uses \texttt{.} (dot), not \texttt{*} (star). 
    The star performs element-wise multiplication.
\end{tipbox}

\subsection{Constructing Special Matrices}

\begin{mathematicabox}
    \begin{lstlisting}[language=Mathematica, basicstyle=\ttfamily\small]
(* Tridiagonal matrix *)
n = 5;
tridiag = Table[
  Which[i==j, 2, i==j-1, 4, i==j+1, 3, True, 0],
  {i,n}, {j,n}
];

(* Using DiagonalMatrix *)
tridiag = DiagonalMatrix[ConstantArray[2,n]] + 
          DiagonalMatrix[ConstantArray[4,n-1], 1] + 
          DiagonalMatrix[ConstantArray[3,n-1], -1];

(* Submatrix: rows {1,3,4}, columns {3,4} *)
sub = tridiag[[{1,3,4}, {3,4}]];
    \end{lstlisting}
\end{mathematicabox}

\subsection{Solving Linear Systems}

\begin{definitionbox}[Methods for Solving \(A\mathbf{x} = \mathbf{b}\)]
    \begin{enumerate}
        \item \textbf{Inverse method:} \(\mathbf{x} = A^{-1}\mathbf{b}\)
        \item \textbf{LinearSolve:} \texttt{LinearSolve[A, b]}
        \item \textbf{Cramer's Rule:} \(x_i = \frac{\det(A_i)}{\det(A)}\)
        \item \textbf{Row Reduction:} \texttt{RowReduce[Join[A, Transpose[\{b\}], 2]]}
        \item \textbf{LU Decomposition:} \(A = LU\)
    \end{enumerate}
\end{definitionbox}

\begin{mathematicabox}
    \begin{lstlisting}[language=Mathematica, basicstyle=\ttfamily\small]
A = {{1,1,2},{2,4,-3},{3,6,-5}};
b = {9,1,0};

(* Method 1: Inverse *)
x = Inverse[A] . b

(* Method 2: LinearSolve (recommended) *)
x = LinearSolve[A, b]

(* Method 3: Cramer's Rule *)
detA = Det[A];
x1 = Det[ReplacePart[A, 1 -> b]]/detA;

(* Method 4: LU Decomposition *)
{lu, p, c} = LUDecomposition[A];
x = LinearSolve[lu, b, p, c];
    \end{lstlisting}
\end{mathematicabox}

\subsection{Network Flow Problems}

\begin{definitionbox}[Kirchhoff's Current Law]
    At each junction: \textbf{Flow in = Flow out}
    
    This creates a system of linear equations.
\end{definitionbox}

% ============================================================
% SECTION 5: ASSIGNMENT 4B - LINEAR INDEPENDENCE & EIGENVALUES
% ============================================================
\section{Assignment 4B: Linear Independence \& Eigenvalues}

\begin{definitionbox}[Learning Goals]
    Test linear independence, find bases/coordinates, apply rank-nullity,
    compute eigenvalues/eigenvectors, and verify diagonalization in Mathematica 13.
\end{definitionbox}

\subsection{Linear Independence}

\begin{definitionbox}[Linear Independence]
    Vectors \(\mathbf{v}_1, \mathbf{v}_2, \ldots, \mathbf{v}_n\) are linearly independent if:
    \[ c_1\mathbf{v}_1 + c_2\mathbf{v}_2 + \cdots + c_n\mathbf{v}_n = \mathbf{0} \implies c_1 = c_2 = \cdots = c_n = 0 \]
    
    \textbf{Test:} Form matrix with vectors as columns. If \(\det \neq 0\) (or rank = n), they are independent.
\end{definitionbox}

\begin{mathematicabox}
    \begin{lstlisting}[language=Mathematica, basicstyle=\ttfamily\small]
(* Test linear independence *)
vectors = {{1,-1,2}, {1,-2,1}, {1,1,4}};
matrix = Transpose[vectors];
rank = MatrixRank[matrix];
isIndependent = (rank == 3)

(* Alternative: Check determinant *)
Det[matrix] != 0

(* For polynomials: convert to coefficient vectors *)
p1 = {0, -1, 0, 1};  (* x^3 - x *)
p2 = {4, 0, 2, 0};   (* 2x^2 + 4 *)
p3 = {6, 2, 3, -2};  (* -2x^3+3x^2+2x+6 *)
    \end{lstlisting}
\end{mathematicabox}

\subsection{Basis and Coordinate Vectors}

\begin{definitionbox}[Coordinate Vector]
    The coordinate vector of \(\mathbf{v}\) relative to basis \(S = \{\mathbf{v}_1, \ldots, \mathbf{v}_n\}\) satisfies:
    \[ \mathbf{v} = c_1\mathbf{v}_1 + \cdots + c_n\mathbf{v}_n \]
    
    Solve: \([\mathbf{v}_1 \ \mathbf{v}_2 \ \cdots \ \mathbf{v}_n][c] = [\mathbf{v}]\)
\end{definitionbox}

\begin{mathematicabox}
    \begin{lstlisting}[language=Mathematica, basicstyle=\ttfamily\small]
(* Find coordinate vector *)
basis = {{1,0,0}, {2,2,0}, {3,3,3}};
v = {2,-1,3};
coeff = LinearSolve[Transpose[basis], v]

(* For polynomials *)
p1 = {1,1,0}; p2 = {1,0,1}; p3 = {0,1,1};
p = {2,-1,1};
coeff = LinearSolve[Transpose[{p1,p2,p3}], p]
    \end{lstlisting}
\end{mathematicabox}

\subsection{Rank and Nullity}

\begin{definitionbox}[Dimension Theorem]
    \[ \rank(A) + \nullity(A) = n \]
    
    where \(n\) is the number of columns of \(A\).
    
    \begin{itemize}
        \item \textbf{Rank:} number of linearly independent rows/columns
        \item \textbf{Nullity:} dimension of null space = number of free variables
    \end{itemize}
\end{definitionbox}

\begin{mathematicabox}
    \begin{lstlisting}[language=Mathematica, basicstyle=\ttfamily\small]
A = {{179, 676, 785, 771},
     {428, 361, 283, 917},
     {640, 512, 154, 303}};
rank = MatrixRank[A];
nullity = Dimensions[A][[2]] - rank;

(* Basis for null space *)
nullSpaceBasis = NullSpace[A];
    \end{lstlisting}
\end{mathematicabox}

\subsection{Eigenvalues and Eigenvectors}

\begin{definitionbox}[Eigenvalue Equation]
    \[ A\mathbf{v} = \lambda\mathbf{v} \quad \text{where } \mathbf{v} \neq \mathbf{0} \]
    
    Characteristic polynomial: \(\det(A - \lambda I) = 0\)
\end{definitionbox}

\begin{formulabox}
    \textbf{Properties:}
    \begin{align*}
    \text{Trace:} &\quad \tr(A) = \sum \lambda_i \\
    \text{Determinant:} &\quad \det(A) = \prod \lambda_i \\
    \text{Algebraic multiplicity:} &\quad \text{multiplicity as root of char poly} \\
    \text{Geometric multiplicity:} &\quad \dim(\text{null space of } A - \lambda I)
    \end{align*}
\end{formulabox}

\begin{mathematicabox}
    \begin{lstlisting}[language=Mathematica, basicstyle=\ttfamily\small]
A = {{4,0,1}, {2,3,2}, {1,0,4}};

(* Characteristic polynomial *)
charPoly = CharacteristicPolynomial[A, lam]

(* Eigenvalues and eigenvectors *)
eigenvalues = Eigenvalues[A]
eigenvectors = Eigenvectors[A]
{values, vectors} = Eigensystem[A]

(* Eigenspace basis for specific lam *)
lam = 3;
eigenspace = NullSpace[A - lam*IdentityMatrix[3]]
    \end{lstlisting}
\end{mathematicabox}

\subsection{Diagonalization}

\begin{definitionbox}[Diagonalization]
    A matrix \(A\) is diagonalizable if there exists invertible \(P\) such that:
    \[ P^{-1}AP = D \]
    
    where \(D\) is diagonal with eigenvalues.
    
    \textbf{Condition:} geometric multiplicity = algebraic multiplicity for every eigenvalue.
\end{definitionbox}

\begin{mathematicabox}
    \begin{lstlisting}[language=Mathematica, basicstyle=\ttfamily\small]
(* Check diagonalizable *)
{vals, vecs} = Eigensystem[A];
isDiagonalizable = (MatrixRank[Transpose[vecs]] == Length[A])

(* Construct P and D *)
P = Transpose[vecs];
Dmat = DiagonalMatrix[vals];

(* Verify P^{-1}AP = D *)
Simplify[Inverse[P] . A . P] == Dmat

(* Cayley-Hamilton Theorem *)
charPoly = CharacteristicPolynomial[A, lam];
charPolyA = charPoly /. lam -> A;
Simplify[charPolyA]  (* Should give zero matrix *)
    \end{lstlisting}
\end{mathematicabox}

% ============================================================
% MASTER FORMULA SHEET
% ============================================================
\section{Master Formula Sheet}

\begin{center}
    \begin{tabular}{|c|c|}
        \hline
        \multicolumn{2}{|c|}{\textbf{\color{primary} CALCULUS FORMULAS}} \\
        \hline
        Definite Integral & \(\int_a^b f(x)\,dx = F(b)-F(a)\) \\
        Average Value & \(f_{\text{avg}} = \frac{1}{b-a}\int_a^b f(x)\,dx\) \\
        Area Between Curves & \(\int_a^b |f(x)-g(x)|\,dx\) \\
        Volume (Disk) & \(\pi\int_a^b [R(x)]^2\,dx\) \\
        Volume (Shell) & \(2\pi\int_a^b x f(x)\,dx\) \\
        Arc Length & \(\int_a^b \sqrt{1+[f'(x)]^2}\,dx\) \\
        Surface Area & \(2\pi\int_a^b f(x)\sqrt{1+[f'(x)]^2}\,dx\) \\
        \hline
        \multicolumn{2}{|c|}{\textbf{\color{primary} LINEAR ALGEBRA FORMULAS}} \\
        \hline
        Distance: Point to Line & \(\frac{|(\mathbf{P-Q})\times\mathbf{v}|}{|\mathbf{v}|}\) \\
        Distance: Point to Plane & \(\frac{|Ax_0+By_0+Cz_0+D|}{\sqrt{A^2+B^2+C^2}}\) \\
        Angle Between Lines & \(\cos\theta = \frac{|\mathbf{v}_1\cdot\mathbf{v}_2|}{|\mathbf{v}_1||\mathbf{v}_2|}\) \\
        Area of Triangle & \(\frac{1}{2}|\mathbf{AB}\times\mathbf{AC}|\) \\
        Volume (Parallelepiped) & \(|\mathbf{a}\cdot(\mathbf{b}\times\mathbf{c})|\) \\
        Rank-Nullity & \(\rank(A) + \nullity(A) = n\) \\
        Characteristic Equation & \(\det(A-\lambda I)=0\) \\
        Cayley-Hamilton & \(p(A)=0\) \\
        Diagonalization & \(A = PDP^{-1}\) \\
        \hline
        \multicolumn{2}{|c|}{\textbf{\color{primary} MEANS FORMULAS}} \\
        \hline
        Arithmetic Mean & \(\frac{1}{n}\sum_{i=1}^n x_i\) \\
        Geometric Mean & \((\prod_{i=1}^n x_i)^{1/n}\) \\
        Harmonic Mean & \(\frac{n}{\sum_{i=1}^n 1/x_i}\) \\
        \hline
    \end{tabular}
\end{center}

% ============================================================
% MATHEMATICA QUICK REFERENCE
% ============================================================
\section{Wolfram Mathematica 13 Quick Reference}

\begin{center}
    \begin{tabular}{|l|l|}
        \hline
        \textbf{Task} & \textbf{Command} \\
        \hline
        Define vector & \texttt{v = \{a,b,c\}} \\
        Dot product & \texttt{v1 . v2} \\
        Cross product & \texttt{Cross[v1, v2]} \\
        Magnitude & \texttt{Norm[v]} \\
        Unit vector & \texttt{v/Norm[v]} \\
        Matrix multiplication & \texttt{A . B} \\
        Determinant & \texttt{Det[A]} \\
        Inverse & \texttt{Inverse[A]} \\
        Transpose & \texttt{Transpose[A]} \\
        Trace & \texttt{Tr[A]} \\
        Rank & \texttt{MatrixRank[A]} \\
        Null space & \texttt{NullSpace[A]} \\
        Eigenvalues & \texttt{Eigenvalues[A]} \\
        Eigenvectors & \texttt{Eigenvectors[A]} \\
        Solve system & \texttt{LinearSolve[A, b]} \\
        LU decomposition & \texttt{LUDecomposition[A]} \\
        Definite integral & \texttt{Integrate[f[x], \{x,a,b\}]} \\
        Numerical integral & \texttt{NIntegrate[f[x], \{x,a,b\}]} \\
        Summation & \texttt{Sum[expr, \{i,imin,imax\}]} \\
        Do loop & \texttt{Do[expr, \{i,imin,imax\}]} \\
        For loop & \texttt{For[start, test, inc, body]} \\
        Random numbers & \texttt{RandomInteger[\{min,max\}, n]} \\
        \hline
    \end{tabular}
\end{center}

% ============================================================
% EXAM TIPS
% ============================================================
\section{Exam Tips \& Common Pitfalls}

\begin{tipbox}
    \textbf{General Tips:}
    \begin{itemize}
        \item Always check dimensions before matrix multiplication
        \item For area between curves, identify which function is on top
        \item When using Cramer's rule, ensure \(\det(A) \neq 0\)
        \item For geometric means, all numbers must be positive
        \item In free fall problems, remember \(g = 32\ \text{ft/s}^2\) or \(9.8\ \text{m/s}^2\)
    \end{itemize}
\end{tipbox}

\begin{tipbox}
    \textbf{Common Mistakes to Avoid:}
    \begin{itemize}
        \item Matrix multiplication uses \texttt{.} not \texttt{*}
        \item Don't forget the \(2\pi\) in surface area formulas
        \item For volume by shells, radius is distance to axis of revolution
        \item In Descartes' rule, ignore zero coefficients when counting sign changes
        \item Geometric multiplicity = dimension of eigenspace, not number of eigenvectors
    \end{itemize}
\end{tipbox}

% ============================================================
% END OF DOCUMENT
% ============================================================
\end{document}
